% Fichier produit mécaniquement depuis la série « planification » du registre.
% Ne pas modifier à la main : docs/planification/to_tex.py le réécrit.

\begin{tikzpicture}[x=1.1cm, y=-0.5cm, font=\footnotesize]
\node[anchor=south, font=\footnotesize] at (0.5, 0) {S30};
\draw[griscartouche, line width=0.3pt] (0, 0) -- (0, 8);
\node[anchor=south, font=\footnotesize] at (1.5, 0) {S31};
\draw[griscartouche, line width=0.3pt] (1, 0) -- (1, 8);
\node[anchor=south, font=\footnotesize] at (2.5, 0) {S32};
\draw[griscartouche, line width=0.3pt] (2, 0) -- (2, 8);
\node[anchor=south, font=\footnotesize] at (3.5, 0) {S33};
\draw[griscartouche, line width=0.3pt] (3, 0) -- (3, 8);
\node[anchor=south, font=\footnotesize] at (4.5, 0) {S34};
\draw[griscartouche, line width=0.3pt] (4, 0) -- (4, 8);
\node[anchor=south, font=\footnotesize] at (5.5, 0) {S35};
\draw[griscartouche, line width=0.3pt] (5, 0) -- (5, 8);
\node[anchor=south, font=\footnotesize] at (6.5, 0) {S36};
\draw[griscartouche, line width=0.3pt] (6, 0) -- (6, 8);
\draw[griscartouche, line width=0.3pt] (7, 0) -- (7, 8);
\node[anchor=east, font=\footnotesize\itshape] at (-0.1, 0.5) {Observation au poste};
\node[diamond, fill=accentgarde, inner sep=2.4pt] at (1.5, 0.5) {};
\node[anchor=east] at (-0.1, 1.5) {Comprendre le service et relever les écrans};
\fill[accentgarde!75] (0, 1.2) rectangle (5, 1.8);
\node[anchor=east] at (-0.1, 2.5) {Établir les sources et les chiffres de référence};
\fill[accentgarde!75] (2, 2.2) rectangle (5, 2.8);
\node[anchor=east] at (-0.1, 3.5) {Concevoir et générer le jeu de données};
\fill[accentgarde!75] (2, 3.2) rectangle (5, 3.8);
\node[anchor=east] at (-0.1, 4.5) {Construire la chaîne de traitement};
\fill[accentgarde!75] (2, 4.2) rectangle (5, 4.8);
\node[anchor=east] at (-0.1, 5.5) {Rapprocher les identités et mesurer la qualité};
\fill[accentgarde!75] (3, 5.2) rectangle (4, 5.8);
\node[anchor=east] at (-0.1, 6.5) {Restituer};
\fill[accentgarde!75] (2, 6.2) rectangle (5, 6.8);
\node[anchor=east] at (-0.1, 7.5) {Rédiger};
\fill[accentgarde!75] (2, 7.2) rectangle (7, 7.8);
\end{tikzpicture}
